\documentclass[10pt,a4paper]{amsart}
\usepackage[margin=19mm]{geometry}
\usepackage{amsmath,amssymb,mathtools}
\usepackage[scr=rsfs,cal=euler]{mathalfa}
\usepackage{xfrac}
\usepackage[unicode,hidelinks,bookmarks=false]{hyperref}
\newcommand{\ie}{i.e.\ }
\newcommand{\cf}{cf.\ }
\newcommand{\resp}{respectively\ }
\newcommand{\Subsection}{\subsection}
\providecommand{\nobreakdash}{\nobreak\hbox{-}}
\setlength{\emergencystretch}{2em}
\multlinegap0pt
\usepackage{enumerate}
\usepackage{dsfont}
\usepackage{mathtools}
\usepackage{cancel}
\usepackage{tikz}
\usepackage{amssymb}
\usepackage[shortlabels]{enumitem}
\newtheorem{proposition}{Proposition}
\newcommand{\tr}{\mbox{tr}}
\newcommand{\ubar}{\underline{u}}
\title{Dynamical Formation of Black Holes due to Boundary Effect in Vacuum}
\author{Puskar Mondal, Shing-Tung Yau}
\date{Source: arXiv:2511.09508v5 (CC BY 4.0)}
\begin{document}
\maketitle

\noindent\emph{Source and licence.} Dynamical Formation of Black Holes due to Boundary Effect in Vacuum. Version arXiv:2511.09508v5 (CC BY 4.0), distributed by arXiv under the Creative Commons Attribution 4.0 International licence, http://creativecommons.org/licenses/by/4.0/, as recorded in the arXiv licence metadata for this exact version and shown on the abstract page https://arxiv.org/abs/2511.09508v5. Full licence notice: \texttt{licenses/arxiv-2511.09508-CC-BY-4.0.txt}.\par\smallskip

\noindent\textit{Editorial scope.} Counted unit: the article's proposition prop:isoperimetric-uniform (source0.tex lines 2452--2470). The article's complete original proof of it is delivered with it (source0.tex lines 2471--2587, one \texttt{proof} environment of 117 lines). No mathematics is written, repaired or re-derived: the statement and the proof are the article's own lines, in the article's own notation, and no retained line is substituted or re-flowed.  No further source-local context is needed: every cross-reference of every retained line resolves inside the two delivered spans.  Nothing was omitted from any retained span, so the package carries no omission record.  No retained line contains a citation, so the wrapper carries no bibliography and no bibliography entry.  No retained line contains a figure environment, an image-inclusion command or drawing code, so no image is delivered and the package declares no asset dependency.  Editorial formatting only, in addition: an AMS \texttt{amsart} preamble that restates the article's own declarations and macros used by the retained lines, namely the environments \texttt{proposition} on separate counters, together with the standard packages those lines need.  The retained environments print in this document as Proposition 1 in order of appearance, whereas the article prints the same items with the numbering of its own full document.  The complete native source of the article as distributed in the arXiv e-print for this exact version is bundled unchanged as \texttt{sources/188986/source0.tex}.
\begin{proposition}[Uniform control of the isoperimetric constant of $S_{u,\ubar}$]
\label{prop:isoperimetric-uniform}
Assume the initial data hypotheses and the bootstrap bounds \textnormal{(2.10)} for the double--null development. Then for every sphere of the foliation
\[
S_{u,\ubar},\qquad u\in [u_{\infty},-a],\ \ubar\in[0,1],
\]
the isoperimetric constant satisfies the uniform estimate
\begin{equation}\label{eq:isoperimetric-bound}
I(S_{u,\ubar})\;\le\;\frac{1}{\pi}.
\end{equation}
Here $I(S,\gamma)$ denotes the isoperimetric constant of the Riemannian $2$–sphere $(S,\gamma)$, defined by
\[
I(S,\gamma)
:=
\sup_{U\subset S}
\frac{\min\{\text{Area}(U),\text{Area}(U^{c})\}}{\text{Per}(\partial U)^2},
\]
where the supremum ranges over all domains $U\subset S$ with smooth boundary.
\end{proposition}

\begin{proof}
We first note that the isoperimetric constant is scale invariant: if $\gamma'=\lambda^{2}\gamma$ on a surface $S$, then
\[
\frac{\min\{\text{Area}_{\gamma'}(U),\text{Area}_{\gamma'}(U^{c})\}}{\text{Per}_{\gamma'}(\partial U)^2}
=
\frac{\lambda^{2}\min\{\text{Area}_{\gamma}(U),\text{Area}_{\gamma}(U^{c})\}}{\lambda^{2}\text{Per}_{\gamma}(\partial U)^2},
\]
hence $I(S,\gamma')=I(S,\gamma)$. Therefore it suffices to prove the estimate for the renormalized metric
\[
\tilde\gamma_{u,\ubar}:=|u|^{-2}\gamma_{u,\ubar}
\quad\text{on }S_{u,\ubar}.
\]

\medskip
\noindent
Suppose a metric $\gamma$ on $S^2$ satisfies, for some $\varepsilon\in(0,1)$,
\begin{equation}\label{eq:bilip-proof}
(1-\varepsilon)\gamma_{0}\le \gamma \le (1+\varepsilon)\gamma_{0}
\end{equation}
as quadratic forms, where $\gamma_{0}$ is the unit round metric. Then for every smooth domain $U\subset S$,
\[
\text{Area}_{\gamma}(U)\le (1+\varepsilon)\text{Area}_{\gamma_{0}}(U), 
\qquad
\text{Area}_{\gamma}(U^{c})\le (1+\varepsilon)\text{Area}_{\gamma_{0}}(U^{c}),
\]
and
\[
\text{Per}_{\gamma}(\partial U)\ge (1-\varepsilon)^{1/2}\text{Per}_{\gamma_{0}}(\partial U),
\quad\Rightarrow\quad
\text{Per}_{\gamma}(\partial U)^2\ge (1-\varepsilon)\text{Per}_{\gamma_{0}}(\partial U)^2.
\]
Hence
\[
\frac{\min\{\text{Area}_{\gamma}(U),\text{Area}_{\gamma}(U^{c})\}}{\text{Per}_{\gamma}(\partial U)^2}
\le
\frac{1+\varepsilon}{1-\varepsilon}
\frac{\min\{\text{Area}_{\gamma_{0}}(U),\text{Area}_{\gamma_{0}}(U^{c})\}}{\text{Per}_{\gamma_{0}}(\partial U)^2}.
\]
Taking the supremum gives
\begin{equation}\label{eq:I-bilip-proof}
I(S,\gamma)\le \frac{1+\varepsilon}{1-\varepsilon} I(S,\gamma_{0}).
\end{equation}
For the unit round sphere one has $I(S^{2},\gamma_{0})=\frac{1}{2\pi}$. Thus if $\varepsilon\le \frac13$, then
\begin{equation}\label{eq:I-target-proof}
I(S,\gamma)\le \frac{1+\varepsilon}{1-\varepsilon}\frac{1}{2\pi}\le \frac{1}{\pi}.
\end{equation}

\medskip
\noindent
In the double–null foliation the induced metric satisfies the transport equations
\[
\nabla_{4}\gamma_{AB}=2\chi_{AB},
\qquad
\nabla_{3}\gamma_{AB}=2\underline{\chi}_{AB},
\]
that is,
\[
\nabla_{4}\gamma_{AB}=(\tr\chi)\gamma_{AB}+2\hat\chi_{AB},
\qquad
\nabla_{3}\gamma_{AB}=(\tr\underline{\chi})\gamma_{AB}+2\underline{\hat\chi}_{AB}.
\]
Define $\tilde\gamma_{AB}=|u|^{-2}\gamma_{AB}$. Using $\nabla_{4}u=0$ and the standard normalization of $\nabla_{3}u$, one obtains schematically
\begin{align}
\nabla_{4}\tilde\gamma_{AB}
&=
\Big(\tr\chi-\frac{2}{|u|}\Big)\tilde\gamma_{AB}
+2|u|^{-2}\hat\chi_{AB},
\label{eq:tildegamma4-proof}
\\
\nabla_{3}\tilde\gamma_{AB}
&=
\Big(\tr\underline{\chi}+\frac{2}{|u|}\Big)\tilde\gamma_{AB}
+2|u|^{-2}\underline{\hat\chi}_{AB}
+\text{l.o.t.}
\label{eq:tildegamma3-proof}
\end{align}
By the initial data assumptions on $H_{u_\infty}$ we have
\[
\tilde\gamma(u_\infty,0)=\gamma_{0}+O(a^{-1/2}|u_\infty|^{-3}).
\]
The bootstrap bounds (2.10) give, uniformly for $u\in[u_\infty,-a]$, $\ubar\in[0,1]$,
\[
\Big|\tr\chi-\frac{2}{|u|}\Big|\lesssim a^{-1/2}|u|^{-1},
\qquad
|\hat\chi|\lesssim a^{-1/2}|u|^{-1},
\]
\[
\Big|\tr\underline{\chi}+\frac{2}{|u|}\Big|\lesssim |u|^{-2},
\qquad
|\underline{\hat\chi}|\lesssim a^{1/2}|u|^{-2}.
\]
Since $\ubar\in[0,1]$ and $|u|\ge a\gg1$, Grönwall’s inequality applied to
\eqref{eq:tildegamma4-proof} along the $\nabla_{4}$ direction and to
\eqref{eq:tildegamma3-proof} along the $\nabla_{3}$ direction yields
\begin{equation}\label{eq:C0-close-proof}
\sup_{u\in[u_\infty,-a],\ \ubar\in[0,1]}
\|\tilde\gamma_{u,\ubar}-\gamma_{0}\|_{C^{0}(S_{u,\ubar})}
\lesssim a^{-1/2}.
\end{equation}
For $a\ge a_{0}$ sufficiently large, this implies the bilipschitz comparison
\eqref{eq:bilip-proof} with some $\varepsilon\le \frac13$, uniformly for all
$S_{u,\ubar}$ in the slab.

\medskip
\noindent
By scale invariance,
\[
I(S_{u,\ubar},\gamma_{u,\ubar})
=
I(S_{u,\ubar},\tilde\gamma_{u,\ubar}).
\]
Combining \eqref{eq:C0-close-proof} with \eqref{eq:I-target-proof} gives
\[
I(S_{u,\ubar})\le \frac{1}{\pi}
\]
uniformly for all $u\in[u_\infty,-a]$, $\ubar\in[0,1]$. This proves the proposition.
\end{proof}

\end{document}
